\newcommand{\panelheight}{0.14\textheight} % adjust as needed

\begin{figure}[t]
    \centering
    \caption{OSS Development Workflow and Social Network}
    % ---------- Subfigure 1 (caption on TOP) ----------
    \begin{subfigure}[t]{0.48\textwidth}
        \centering
        \caption{Release Production Workflow}
        \label{fig:development_pipeline}
        % force TikZ to match the same height as the PNG
        \vspace{0pt}
        \resizebox{!}{\panelheight}{%
            % issues_prs_panel.tikz
\def\colx{0cm}      % left edge of column 1
\def\colsep{2.8cm}    % horizontal gap between columns (reduce for shorter arrows)
\def\rowsep{7mm}   % vertical gap between rows (increase if still overlapping)
\def\titleoffset{4mm}   % <--- adjust this number to control the gap

\begin{tikzpicture}[
  % --- node styles: small inner sep so boxes are barely bigger than their text ---
  issue/.style={draw=black, rounded corners, thick, inner sep=2pt, font=\small},
  pr/.style={draw=black, rounded corners, thick, inner sep=2pt, font=\small\bfseries, fill=gray!10},
  reviewApproved/.style={draw=black, rounded corners, thick, inner sep=2pt, font=\footnotesize\bfseries, fill=green!12},
  reviewNot/.style={draw=black, rounded corners, thick, inner sep=2pt, font=\footnotesize\bfseries, fill=red!8},
  % make release a visible circle; minimum width forces round shape, align=center allows linebreaks
  release/.style={draw=black, rounded corners,  thick, minimum width=14mm, inner sep=3pt, align=center, font=\small\bfseries, fill=black!6},
  % --- arrow style: slightly larger stealth tip so it's clearly visible ---
  arrow/.style={-{Stealth[length=5pt,width=5pt]}, line width=0.9pt, shorten >=1pt, shorten <=1pt},
  columnTitle/.style={anchor=north, font=\large}
]

% Column titles
\node[columnTitle] (Cissues)  at (\colx,0+\titleoffset) {Issues};
\node[columnTitle] (Cprs)     at (\colx+\colsep,0+\titleoffset) {PRs};
\node[columnTitle] (Creview)  at (\colx+2*\colsep,0+\titleoffset) {Review Decision};
\node[columnTitle] (Crelease) at (\colx+3*\colsep,0+\titleoffset) {Release};

% ---------- rows (6 issues; 4 PRs) ----------
% Use exact multiples of \rowsep to avoid accidental overlap
\node[issue] (Issue1) at (\colx,-1*\rowsep) {Issue 1};
\node[issue] (Issue2) at (\colx,-2*\rowsep) {Issue 2};
\node[issue] (Issue3) at (\colx,-3*\rowsep) {Issue 3};
\node[issue] (Issue4) at (\colx,-4*\rowsep) {Issue 4};
\node[issue] (Issue5) at (\colx,-5*\rowsep) {Issue 5}; % disconnected
\node[issue] (Issue6) at (\colx,-6*\rowsep) {Issue 6}; % disconnected

% PR column (one PR per linked issue)
\node[pr] (PRA) at (\colx+\colsep,-1*\rowsep) {PR A};
\node[pr] (PRB) at (\colx+\colsep,-2*\rowsep) {PR B};
\node[pr] (PRC) at (\colx+\colsep,-3*\rowsep) {PR C};
\node[pr] (PRD) at (\colx+\colsep,-4*\rowsep) {PR D}; % not approved

% Review Decision column (three Approved, one Not approved)
\node[reviewApproved] (RevA) at (\colx+2*\colsep,-1*\rowsep) {Approved};
\node[reviewApproved] (RevB) at (\colx+2*\colsep,-2*\rowsep) {Approved};
\node[reviewApproved] (RevC) at (\colx+2*\colsep,-3*\rowsep) {Approved};
\node[reviewNot]      (RevD) at (\colx+2*\colsep,-4*\rowsep) {Not approved};

% Release (place slightly between rows 2 and 3 vertically)
% --- Use \\ to force a line break and ensure the text wraps to two lines inside the circle ---
\node[release] (Rel) at (\colx+3*\colsep,-2.25*\rowsep) {Software\\Release};

% ---------- connections (compact arrows) ----------
% Issue -> PR (one-to-one)
\draw[arrow] (Issue1) -- (PRA);
\draw[arrow] (Issue2) -- (PRB);
\draw[arrow] (Issue3) -- (PRC);
\draw[arrow] (Issue4) -- (PRD);

% PR -> Review Decision
% using straight short lines because columns are close (colsep small)
\draw[arrow] (PRA) -- (RevA);
\draw[arrow] (PRB) -- (RevB);
\draw[arrow] (PRC) -- (RevC);
\draw[arrow] (PRD) -- (RevD);

% Approved -> Release (only for approved review decisions)
% small bends so arrows are visually separate and short
\draw[arrow] (RevA) to[bend right=8] (Rel);
\draw[arrow] (RevB) -- (Rel);
\draw[arrow] (RevC) to[bend left=8] (Rel);
% RevD (Not approved) has no arrow to release

\end{tikzpicture}

        }
    \end{subfigure}
    \hfill
    % ---------- Subfigure 2 (caption on TOP) ----------
    \begin{subfigure}[t]{0.48\textwidth}
        \centering
        \caption{Social Network for \texttt{boto/boto3}}
        \label{fig:boto_social_network}
        \vspace{0pt}
        \includegraphics[height=\panelheight]{source/raw/gephi/boto_2020_1.png}
    \end{subfigure}

    \vspace{4mm} % optional spacing between plots and notes

    % ---------- Figure Notes ----------
    \begin{minipage}{1\textwidth}
        \footnotesize
        \textbf{Figure notes:} 
        Panel~\ref{fig:development_pipeline} depicts the development process for an OSS to produce a new software release. 
        Panel~\ref{fig:boto_social_network} depicts the social network for \texttt{boto/boto3}, the Amazon Web Services Software Development Kit for Python, during the first half of 2020. 
        Each node represents a member and is scaled by its degree, defined as the number of members they interact with.
        Each edge between nodes represents interactions between members; edge thickness is positively related to the quantity of interactions between members. 
        The three nodes representing the three \textit{single-period key members} and their corresponding edges are colored in red.
    \end{minipage}

\end{figure}
